\setcounter{ex}{0} %Trần Quốc Toản-TĐ
\subsection{Trắc nghiệm (2 điểm)}
\Opensolutionfile{ans}[ans/ans10] 

%Câu 1
\begin{ex}
	Tìm điều kiện xác định của biểu thức $\dfrac{1}{\sqrt{x-2}}+\sqrt{5-x}$.
	\choice
	{$2\leq x<5$}
	{$x>2$}
	{\True $2<x\leq 5$}
	{$2\leq x\leq 5$}
\end{ex}

%Câu 2
\begin{ex}
	Hệ nào sau đây \textbf{không} phải là hệ phương trình bậc nhất hai ẩn?
	\choice
	{\True $\heva{&-x+3 y=2 \\ &0 x+0 y=-1}$}
	{$\heva{&3 x-y=3 \\ &2 x+y=-8}$}
	{$\heva{&-5 x+y=4 \\ &2 x+3 y=0}$}
	{$\heva{&2 x-5 y=7 \\& \dfrac{1}{2} x-y=1}$}
\end{ex}

%Câu 3
\begin{ex}
Nghiệm của phương trình $(4x^2+4x+1)-(2x+3)^2=0$ là:
\choice
{\True $x=-1$}
{$x=0$}
{$x=3$}
{Vô nghiệm}
\end{ex}

%Câu 4
\begin{ex}
Cho $ 2 $ số thực $a, b$ thỏa mãn $a<b-2$. Tìm khẳng định đúng.
\choice
{\True $-2a>-2b+4$}
{$2-b>-a$}
{$a+2 > b$}
{$a-b>2$}
\end{ex}

%Câu 5
\begin{ex}
Giá trị $x=2$ là một nghiệm của bất phương trình nào dưới đây?
\choice
{$3 x-5<1$}
{$2 x-1 \leq 1$}
{\True $-2 x+7>2$}
{$2 x+1 \leq-4$}
\end{ex}

%Câu 6
\begin{ex}
Cho $\triangle DEF$ vuông tại $F$. Khi đó $\cot \widehat{EDF}$ bằng
\choice
{$\dfrac{FE}{ED}$}
{$\dfrac{FD}{ED}$}
{\True $\dfrac{FD}{FE}$}
{$\dfrac{FE}{FD}$}
\end{ex}

%Câu 7
\begin{ex}
Cho tam giác $ABC$ nhọn nội tiếp đường tròn $(O)$ với $\widehat{BAC}=60^\circ$. Dựng đường kính $AD$ của $(O)$. Tính $\widehat{BDC}$.
\choice
{$60^\circ$}
{$90^\circ$}
{$45^\circ$}
{\True $120^\circ$}
\end{ex}

%Câu 8
\begin{ex}
	Cho tam giác $ABC$ nhọn nội tiếp đường tròn $(O;R)$. Biết rằng $BC=R\sqrt{2}$. Tính $\widehat{BAC}$.
\choice
{$60^\circ$}
{$90^\circ$}
{\True $45^\circ$}
{$120^\circ$}
\end{ex}

%Câu 9
\begin{ex}
Từ một tờ giấy hình tròn bán kính 5 cm, có thể cắt thành hình vuông có diện tích lớn nhất là bao nhiêu?
\choice
{\True 50 $\text{cm}^2$}
{25 $\text{cm}^2$}
{10 $\text{cm}^2$}
{75 $\text{cm}^2$}
\end{ex}

%Câu 10
\begin{ex}
 Một người thợ cần đo độ cao một cây cột điện nhưng không có thước dài. Anh ấy đứng cách gốc cột điện 12 m và dùng thước đo góc đo được góc giữa mặt đất và đỉnh cột điện là $35^\circ$. Hãy tính chiều cao của cột điện (\textit{làm tròn đến chữ số thập phân thứ hai}).
 \choice
 {7 m}
 {\True 8,4 m}
 {9,8 m}
 {12 m}
\end{ex}
    
\Closesolutionfile{ans}
\Opensolutionfile{ansbook}[ans/ansbook10]
\setcounter{bt}{0}
\subsection{Tự luận (8 điểm)}

%%Bài 1:
\begin{bt}[1,5 điểm]
Giải các phương trình và hệ phương trình sau sau:
\begin{listEX}[3]
    \item $(3x-6)(x^2+x+1)=0$
    \item $\dfrac{2}{x-3}+\dfrac{3}{x+3}=\dfrac{3 x-5}{x^2-9}$
    \item  $\begin{cases}4 x+y=7 \\ -x+y=-3\end{cases}$
\end{listEX}
\loigiai{
a) $x=2$ \hspace{2cm} b) $x=-1$ \hspace{2cm} c)  $(2;-1)$
}
\end{bt}

%%Bài 2:
\begin{bt}[1,5 điểm]
Thực hiện phép tính (thu gọn)
\begin{enumerate}
    \item $\sqrt{27}-3 \sqrt{48}+\dfrac{5}{2} \sqrt{108}$
    \item $\dfrac{\sqrt{15}-\sqrt{12}}{\sqrt{5}-2}+\dfrac{6}{\sqrt{3}+1}+\sqrt{(4 \sqrt{3}-12)^2}$
\end{enumerate}
\loigiai{
a) $6\sqrt{3}$ \hspace{2cm} b) 9
}
\end{bt}

%%Bài 3:
\begin{bt}[1,0 điểm]
Như thường ngày, mẹ đưa 115\,000 đồng để nhờ An mua 2 ổ bánh mì và 3 cái bánh bao ở cửa hàng M cho cả nhà ăn sáng là vừa đủ. Hôm nay, cửa hàng M giảm giá $10 \%$ cho tất cả các mặt hàng và An đổi ý muốn ăn bánh bao nên bạn mua 1 ổ bánh mì và 4 cái bánh bao. An tính toán với số tiền mang theo bạn còn dư 7\,000 đồng. Hỏi giá của mỗi ổ bánh mì, mỗi cái bánh bao khi cửa hàng M chưa giảm giá là bao nhiêu?
\loigiai{
Một cái bánh mỳ có giá 20\,000 đồng, một cái bánh bao có giá 25\,000 đồng.
}
\end{bt}

%%Bài 4:
\begin{bt}
[1,0 điểm] \immini{Hải đăng ngày nay là một ngọn tháp (nhà hoặc khung) được thiết kế để chiếu sáng từ một hệ thống đèn và thấu kính với mục đích hỗ trợ cho các hoa tiêu trên biển định hướng và tìm đường.

Một ngọn hải đăng cao $100$ mét được đặt sát bờ biển. Giả sử mặt biển là một cung tròn của đường tròn có tâm là tâm Trái Đất và đỉnh ngọn hải đăng, chân ngọn hải đăng và tâm Trái Đất là ba điểm thẳng hàng. Hỏi khoảng cách giữa vị trí xa nhất các tàu thuyền có thể thấy được ánh sáng từ ngọn hải đăng và chân ngọn hải đăng là bao nhiêu kilomet? Biết rằng bán kính Trái Đất khoảng $6371$ km và góc nâng của đèn tối đa là $90^\circ$. (Làm tròn kết quả đến hàng phần mười)
}{\begin{tikzpicture}
	\path 
	(0:0) coordinate (C)
	(0,1.125) coordinate (A)
	(0,-2.5) coordinate (O)
	($(A)!0.5!(O)$) coordinate (I);
	\path[name path=tron1] (O) circle (2.5);
	\path[name path=tron2] (I) let\p1=($(I)-(O)$) in circle ({veclen(\x1,\y1)});
	\path[name intersections={of=tron1 and tron2,by={B,D}}];
	\draw (A)--(B);
	\draw (O) circle (2.5);
	\begin{scope}[scale=.0625]
		\draw [fill=gray!60!black,smooth,line join=round,gray!60!black] (-6.14475,0)..controls +(60:.5) and +(120:.5)..(-5.727,0.3825)..controls +(60:.5) and +(120:.5)..(-4.1775,0.9825)..controls +(60:.5) and +(120:.5)..(-3.0275,2)..controls +(60:.5) and +(120:.5)..(-2.,2.)--(-1.8575,4.5325)--(-1.6275,3.1825)..controls +(70:1) and +(-110:1)..(1.8,5.5825)--(2.,2.)..controls +(60:.5) and +(120:.5)..(2.9225,1.9325)..controls +(60:.5) and +(120:.5)..(4.2225,1.0825)..controls +(60:.5) and +(120:.5)..(5.9225,0.4325)..controls +(60:.5) and +(120:.5)..(6.7725,0)--cycle;
		\draw [fill=gray!60!black,smooth,gray!60!black] (-1.77,5.8825)--(-1.62,8.7325)..controls +(70:1) and +(-110:1)..(1.48,10.8325)--(1.612,8.0825)..controls +(-110:1) and +(70:1)..(-1.77,5.8825);
		\draw [fill=gray!60!black,gray!60!black] (-1.45,11.4325)--(-1.22,15) arc (180:0:{1.22} and {.5})--(1.22,15)--(1.34,13.3325)..controls +(-110:1) and +(70:1)..(-1.45,11.4325);
		\draw [line width=.1cm,gray!60!black] (-1.18,15)--(-1.22,15.6) arc (180:0:{1.22} and {.5})--(1.22,15.6)--(1.18,15);
		\draw[fill=gray!60!black,gray!60!black] (0,14) ellipse ({2} and {.5});
		\draw[fill=gray!60!black,gray!60!black] (0,17.5) ellipse ({.2} and {.8});
		\draw [fill=gray!60!black,gray!60!black] (0,16.4) circle (.35);
		\draw[fill=gray!60!black,gray!60!black] (-.9,17) arc (180:0:{.9 and .9});
		\draw [line width=.08cm,gray!60!black] (0.87,15)--(0.87,16);
		\draw [line width=.08cm,gray!60!black] (-0.87,15)--(-0.87,16);
		\draw [line width=.08cm,gray!60!black] (0.33,15)--(0.33,16.1);
		\draw [line width=.08cm,gray!60!black] (-0.33,15)--(-0.33,16.1);
		\draw [line width=.06cm,gray!60!black] (0.81,16)--(0.81,17.2);
		\draw [line width=.06cm,gray!60!black] (-0.81,16)--(-0.81,17.2);
	\end{scope}
	\draw[line width=0.2mm, red] (0,0) arc[start angle=90, end angle=43.8, radius=2.5] --cycle;
	\draw[line width=0.2mm, black] (C)--(B);
	\draw (A)--(O)--(B)--(C);
	\foreach \x/\g in{A/90,B/45,O/-90,C/-135}
	\draw[fill=white] (\x) circle (.05) +(\g:.3) node{$\x$};
\end{tikzpicture}
}
 \loigiai{Khoảng $35{,}7$ km.}
 \end{bt}

%%Bài 5:
\begin{bt}[2,5 điểm]
Cho đường tròn $(O;R)$. Từ điểm S nằm ngoài đường tròn $(O)$, vẽ hai tiếp tuyến $SM$ và $SN$ với đường tròn $(O)$ ($M, N$ là hai tiếp điểm). Gọi $H$ là giao điểm của $SO$ và $MN$ . Kẻ đường kính $NE$ của đường tròn $(O)$, $SE$ cắt đường tròn $(O)$ tại điểm $F$ ($E$ khác $F$).
\begin{enumerate}
    \item Chứng minh: $SO \perp MN$ tại $H$.
    \item Chứng minh: $SM^2=SF\cdot SE$ và $SF\cdot SE=SH\cdot SO$
    \item Gọi $K$ là trung điểm của $EF$ và $I$ là giao điểm của $MN$ với $EF$.
    
    Chứng minh: $SF\cdot SE=SK\cdot SI$ từ đó suy ra $\dfrac{1}{SE}+\dfrac{1}{SF}=\dfrac{2}{SI}$.
\end{enumerate}
\loigiai{\immini
{\begin{enumerate}
		\item $MN$ là đường trung trực của $SO$
		\item $\triangle SNF \backsim \triangle SEN$ suy ra $SN^2=SF\cdot SE$
		
		$\triangle SNH \backsim \triangle SON$ suy ra $SN^2=SH\cdot SO$
		\item $\triangle SHI \backsim \triangle SKO$ suy ra $SK\cdot SI=SH\cdot SO$
		
		Chứng minh $SF+SE=2SK$ 
		
		Suy ra $SI(SF+SE)=2SI\cdot SK=2SE\cdot SF$
\end{enumerate}}
   {\begin{tikzpicture}[scale=1, font=\footnotesize]
    \coordinate (O) at (0,0);
    \def\R{2}
    \draw (O) circle(\R);
    \coordinate (S) at (-5,0);
    \path[name path=circleO] (O) circle(\R);
    \path[name path=circleS] (-2.5,0) circle(2.5);
    \path[name intersections={of=circleS and circleO, by={M,N}}];
    \coordinate (E) at ($-1*(N)$);
    \path[name path=os] (O)--(S);
    \path[name path=mn] (M)--(N);
    \path[name intersections={of=os and mn, by=H}];
    \path[name path=se] (S)--(E);
    \path[name intersections={of=circleO and se, by={F1,F2}}]; 
    \coordinate (F) at (F2);
    \path[name intersections={of=se and mn, by=I}];
    \draw (S)--(M)--(O)--cycle;
    \draw (M)--(N)--(E)--(S);
    \draw (S)--(O);
    \coordinate (K) at ($(F)!0.5!(E)$);
    \draw (O)--(K);
    \foreach \p/\pos in {S/left,M/above,N/below,O/below right,F/below,E/above,K/above,I/below left,H/below left}
      {\draw[fill=white] (\p) circle(1.5pt) node[\pos] {$\p$};
      }
\end{tikzpicture}}
}
\end{bt}

%%Bài 6:
\begin{bt}[0,5 điểm]
Giải bất phương trình sau: 
\[\dfrac{x+2021}{2022}+\dfrac{x+2022}{2023}<\dfrac{x+2023}{2024}+\dfrac{x+2024}{2025}\]
\loigiai{
$x=1$

Gợi ý: cộng hai vế của phương trình với -2: 
\[\dfrac{x-1}{2022}+\dfrac{x-1}{2023}>\dfrac{x-1}{2024}+\dfrac{x-1}{2025}\]
}
\end{bt}

\Closesolutionfile{ansbook}